\answer{ex:22:1}{(a)~$32\%$ và $100\%$. (b)~$R_\uparrow-R_\downarrow\ge2\sqrt{2000\,\text{Hz}/1\,\text{s}}\approx89$~Hz, chỉ $4.5\%$ tần số tổng.}
\answer{ex:22:2}{(a)~$n\ge4(p_1q_1+p_2q_2)/\Delta^2\approx560$. (b)~$\approx56$.}
\answer{ex:22:3}{(a)~Cân bằng chi tiết: $\rho PK$ đối xứng. (b)~Tỷ lệ trượt bằng trung bình điều hòa $1/\langle1/P\rangle=0.18$, so với $0.5$ khi không có phản hồi.}
\answer{ex:22:4}{(a)~Hình bình hành diện tích $4h^2=0.04$ (tính cả giới hạn của $p$ thì bằng số là $0.045$). (b)~$\approx0.18$.}
\answer{ex:22:5}{$0.49$: một phần mẫu nuôi cấy trôi vào vùng gần như luôn trượt và lang thang ở đó. Có giới hạn thì $1.00$: trên một tập bị chặn, mật độ $\propto1/P$ dồn vào vùng hấp thụ.}
\answer{ex:22:6}{(a)~Cần $\ge5$ mức; $8$ điện cực là đủ. (b)~Không: cần $r_{\max}\ge25/T=100$~Hz.}
\answer{ex:22:7}{Bài mở. Thời gian tìm kiếm ngẫu nhiên tăng theo $d$ xấp xỉ như tỷ số thể tích $(L/h)^d$, còn tìm kiếm có dấu của sai số thì tăng tuyến tính theo $d$. Thí nghiệm hoán đổi sẽ phân biệt hai giả thuyết: sau cú trượt là kích thích đoán trước được nhưng mạnh, sau cú trúng là kích thích không đoán trước được nhưng yếu; cần vài chục mẫu nuôi cấy mỗi nhóm.}
